Work adjacent to this system falls into three groups, covering networked health monitoring,
embedded queue management, and low-cost stand-alone vital measurement.
Table~\ref{tab:comparison} positions the present design against a representative sample of
each. The pattern it makes visible is that sensing and queueing have been treated as separate
problems, and that the systems which integrate sensing do so by adding a network.

\begin{table*}[t]
  \caption{Representative prior systems compared against the proposed design.
    ``Network required'' means the normal operating path depends on a Wi-Fi, GSM
    or Internet link, and ``n/r'' marks a property the cited work does not
    report.}
  \label{tab:comparison}
  \footnotesize
  \begin{tabular}{@{}llccccc@{}}
    \toprule
    Work & Platform & Network & Queue & Vital & Patient self- & Local clinician \\
         &          & required & arbitration & sensing & registration & display \\
    \midrule
    Dharmoji et al.~\cite{irjet2020}        & ESP8266 + Arduino        & Yes (Wi-Fi)  & No           & Pulse, temp. & No  & No \\
    Rath et al.~\cite{rath2021}             & Raspberry Pi BSN         & Yes (MQTT/CoAP) & No        & BSN sensors  & No  & No \\
    Talluri et al.~\cite{hrtemp2019}        & Arduino Uno + web app    & Yes (web)    & No           & Pulse, temp. & No  & No \\
    Gummadi Lahari et al.~\cite{eqc2022}    & 8-bit PIC                & No           & FIFO         & No           & No  & Yes (LCD) \\
    Arun and Priyesh~\cite{qgsm2013}        & MCU + GSM module         & Yes (GSM)    & FIFO         & No           & No  & Yes \\
    Chin et al.~\cite{arpn2019}             & Arduino                  & n/r          & Priority     & No           & No  & Yes \\
    \midrule
    \textbf{This work}                      & \textbf{ATmega32}        & \textbf{No}  & \textbf{FCFS} & \textbf{Pulse, temp.} & \textbf{Yes} & \textbf{Yes (dual LCD)} \\
    \bottomrule
  \end{tabular}
\end{table*}


\subsection{Networked Health Monitoring}
The largest group couples edge sensing to a remote back end. Dharmoji et al.~\cite{irjet2020}
pair an ESP8266 with an Arduino to read pulse and temperature and publish them to ThingSpeak,
so readings are viewed through a cloud dashboard rather than at the point of care. Rath et
al.~\cite{rath2021} work at larger scale, proposing a body-sensor-network architecture for
hospital management with real-time monitoring and prescription handling, and evaluating MQTT
against CoAP in NS-3 simulation to choose a transport. Talluri et al.~\cite{hrtemp2019} read
heart rate and body temperature from an Arduino Uno and surface them through a web
application.

The remote-access capability these systems provide is one our design deliberately gives up, and the
trade-off they accept is that the reporting path is also the display path. Rath et al. treat
reliability and latency under hospital-scale traffic as the governing concerns, which is the wrong
framing for a chamber where the physician is three metres from the patient.

\subsection{Embedded Queue and Token Management}
The second group automates service order without clinical sensing. The electronic queue control
system of Gummadi Lahari et al.~\cite{eqc2022} is the closest offline analogue to our queue
component. An 8-bit PIC maintains FIFO order across service counters and drives an LCD showing
counter and token number, entirely in software and with no external infrastructure. Arun and
Priyesh~\cite{qgsm2013} build queue handling around a GSM link, reintroducing a network dependency
and a per-message cost. Chin et al.~\cite{arpn2019} address a different limitation, that strict FIFO
ignores patients who should be prioritised, and propose an Arduino-based manager that automatically
advances handicapped persons.

All three show that queue arbitration is comfortably within reach of an 8-bit part. None forms a
patient record or captures a physiological measurement, so a chamber deploying any of them would
still take vitals during the consultation. The priority mechanism of Chin et al. is a capability our
strict-FCFS design lacks, and Section~\ref{sec:limitations} returns to it.

\subsection{Stand-Alone Vital Measurement}
The third group establishes that useful vital acquisition is achievable with inexpensive analog
parts. The LM35 is the conventional choice for body temperature because its output is linear at
10~mV/$^{\circ}$C and needs no linearisation or external conditioning before reaching a
microcontroller ADC~\cite{ti2022}. Optical pulse measurement is less forgiving.
Photoplethysmography infers a pulse from light absorbed by perfused tissue, and the waveform is
sensitive to probe placement, contact pressure, perfusion and above all to motion, which produces
artefacts overlapping the signal band that simple filtering cannot remove~\cite{allen2007ppg}. That
sensitivity is intrinsic to the measurement principle rather than to low-cost hardware, it sets the
accuracy ceiling reported in Section~\ref{sec:results}, and it is why the patient-side workflow
enforces a settling period. Expectations should be calibrated accordingly, since clinical electronic
thermometers are held to laboratory accuracy limits measured in tenths of a
degree~\cite{iso80601}, several times tighter than an uncalibrated LM35 and a 10-bit ADC deliver.

The gap this work addresses is integration rather than any individual capability. Queue
arbitration, guided data entry and vital acquisition have each been demonstrated on comparable
hardware. What has not is all three composed into one offline workflow on a single 8-bit
microcontroller, with the two interfaces separated so registration and consultation proceed
concurrently.
